
\documentclass[10pt]{article}
\usepackage[margin=0.75in]{geometry}
\usepackage[T1]{fontenc}
\usepackage{lmodern}
\usepackage{amsmath}
\usepackage{amsfonts}
\usepackage{booktabs}
\usepackage{array}
\usepackage{graphicx}
\usepackage{bm}
\usepackage{hyperref}

\setlength{\parindent}{0pt}
\setlength{\parskip}{0.45em}
\renewcommand{\arraystretch}{1.14}

\newcommand{\PN}{\ensuremath{P_N}}
\newcommand{\FPN}{\ensuremath{FP_N}}
\newcommand{\rNN}[1]{\ensuremath{r_{#1}^{\mathrm{NN}}}}
\newcommand{\rAugConst}[1]{\ensuremath{r_{#1}^{\mathrm{const,aug}}}}
\newcommand{\rPaperConst}[1]{\ensuremath{r_{#1}^{\mathrm{const,paper}}}}
\newcommand{\cell}[1]{\(#1\)}
\newcommand{\cellbest}[1]{\(\bm{#1}\)}

\title{Manuscript Add-On Sections for the Updated Learning Workflow}
\author{NN\_FPN repository notes}
\date{\today}

\begin{document}
\maketitle

\section{Updated learning workflow for the filtered \texorpdfstring{$P_N$}{PN} model}
The original filtered \(\PN\) solver is retained, but the learned filter is now
trained and selected by a sweep-driven workflow.  For a moment order \(N\), let
\(u_i^n=(u_{i,0}^n,\ldots,u_{i,N}^n)^\top\) denote the cell-centered moment vector
and let \(\phi_i^n=u_{i,0}^n\) denote the scalar flux.  The neural filter is a
cellwise nonnegative map
\[
  \sigma_{f,i}^n = F_{\theta,N}(x_i^n) \ge 0,
\]
where \(x_i^n\) is the feature vector defined below.  The constant-filter
baselines use a single learned scalar \(c_N\ge 0\), i.e.
\[
  \sigma_{f,i}^n = c_N .
\]
The learned filter strength multiplies the same filter coefficients used in the
filtered \(P_N\) update, so the discretization and evaluation metric remain those
of the manuscript.

\subsection{Training objective}
For a training draw \(\omega\), let \(\phi^{\FPN}_{N,\theta}(x_i,T_\omega;\omega)\)
be the filtered \(P_N\) scalar flux produced by the learned filter and let
\(\phi^{\mathrm{ex}}(x_i,T_\omega;\omega)\) be the high-order reference solution.
The updated neural-network and constant-filter sweep models minimize the
relative final-time scalar-flux loss
\[
  \mathcal{L}_N(\theta)
  = \frac{1}{B}\sum_{b=1}^{B}
    \frac{\frac{1}{n_x}\sum_{i=1}^{n_x}
      \left(\phi^{\FPN}_{N,\theta}(x_i,T_b;\omega_b)
      -\phi^{\mathrm{ex}}(x_i,T_b;\omega_b)\right)^2}
    {\max\!\left(\frac{1}{n_x}\sum_{i=1}^{n_x}
      \left(\phi^{\mathrm{ex}}(x_i,T_b;\omega_b)\right)^2,\varepsilon\right)},
  \qquad \varepsilon=10^{-12} .
\]
No auxiliary moment loss is used in the reported updated runs.  The scalar
metric reported in the result tables is unchanged from the manuscript,
\[
  r_N =
  \frac{\left\|\phi^{\FPN}_N-\phi^{\mathrm{ex}}\right\|_2}
       {\left\|\phi^{P_N}-\phi^{\mathrm{ex}}\right\|_2},
\]
so smaller values indicate stronger error reduction relative to the unfiltered
\(P_N\) solution.

\subsection{Training distributions}
The paper-setting baseline uses the original four-family training distribution,
which combines Gaussian, Heaviside, bump, and discontinuous-source initial/source
families with random homogeneous material coefficients.  The updated
sweep-selected workflow replaces this with an augmented distribution
\(\mathcal{D}_{\mathrm{aug}}\).  A draw
\(\omega\sim\mathcal{D}_{\mathrm{aug}}\) samples the initial condition, source,
and material profile from a mixture that includes Gaussian, discontinuous,
vanishing-cross-section-like, source-driven, and Reeds-like cases.  This changes
the training ensemble only; all reported test cases and the definition of
\(r_N\) are unchanged.

\section{Input feature map replacing the baseline network inputs}
For each cell \(i\), define the discrete transport residual and material/source
terms
\[
  d_i = \frac{(A u)_{i+1/2}-(A u)_{i-1/2}}{\Delta x}, \qquad
  c_i = \sigma_{t,i}u_i, \qquad
  s_i = \sigma_{s,i}u_{i,0}, \qquad
  q_i = Q_i,
\]
where \(d_i,c_i\in\mathbb{R}^{N+1}\) and \(s_i,q_i\in\mathbb{R}\).  For any
cellwise channel tensor \(g_{i,k}\), let
\[
  \mathcal{N}_x(g)_{i,k}
  = \frac{g_{i,k}-\frac{1}{n_x}\sum_{j=1}^{n_x}g_{j,k}}
         {\operatorname{std}_{1\le j\le n_x}(g_{j,k})+10^{-10}}
\]
be the per-sample spatial normalization.  The baseline manuscript-style feature
vector is
\[
  z_i^{\mathrm{base}}
  = \Big[
      \mathcal{N}_x(|d_i|),
      \mathcal{N}_x(|c_i|),
      \mathcal{N}_x(|s_i|),
      \mathcal{N}_x(|q_i|)
    \Big]\in\mathbb{R}^{2N+4}.
\]
The sweep also introduced log-magnitude material and residual features,
\[
  L_{\alpha,[a,b]}(g)=
  \operatorname{clip}_{[a,b]}\!\left(\log\!\left(1+\frac{|g|}{\alpha}\right)\right),
  \qquad
  \sigma_{a,i}=\max(\sigma_{t,i}-\sigma_{s,i},0),
\]
and the material vectors
\[
  m_i^{\log}=\Big[
    L_{\alpha,[a,b]}(\sigma_{s,i}),
    L_{\alpha,[a,b]}(\sigma_{t,i}),
    L_{\alpha,[a,b]}(\sigma_{a,i})
  \Big],
\]
\[
  m_i^{\mathrm{all}}=\Big[
    m_i^{\log},
    \frac{\sigma_{s,i}}{\max(|\sigma_{t,i}|,\varepsilon)},
    \frac{\sigma_{a,i}}{\max(|\sigma_{t,i}|,\varepsilon)}
  \Big].
\]
The selected feature maps are
\begin{align*}
  x_i^{(3)} &= \big[z_i^{\mathrm{base}},\,m_i^{\mathrm{all}}\big],
  &\alpha&=0.1, & [a,b]&=[0,20], & \dim x_i^{(3)}&=15, \\
  x_i^{(7)} &= \big[
      L_{0.1,[0,10]}(d_i),
      L_{0.1,[0,10]}(c_i),
      L_{0.1,[0,10]}(s_i),
      L_{0.1,[0,10]}(q_i),
      \mathcal{N}_x(m_i^{\mathrm{all}})
    \big],
  &&& [a,b]&=[0,10], & \dim x_i^{(7)}&=23, \\
  x_i^{(9)} &= \big[z_i^{\mathrm{base}},\,\mathcal{N}_x(m_i^{\log})\big],
  &\alpha&=1.0, & [a,b]&=[0,40], & \dim x_i^{(9)}&=25 .
\end{align*}
The \(N=3\) and \(N=9\) maps retain the normalized baseline residual/source
features, while the \(N=7\) map uses clipped log-magnitude residual/source
features without residual/source normalization.

\section{Selected training settings}
The hyperparameters below are the settings selected by the sweep and then
retrained for 10 independent replicas.  Each run uses Adam and the final-time
scalar-flux loss above.
\begin{center}
\resizebox{\textwidth}{!}{%
\begin{tabular}{c c c c c c c c c c}
\toprule
\(N\) & hidden & batch & epochs & lr & scheduler & wd NN & wd const & feature map & horizons \\
\midrule
3 & 64 & 32 & 100 & 0.003 & constant & \(10^{-3}\) & \(10^{-4}\) & \(x_i^{(3)}\) & \(\{0.5\}\) \\
7 & 512 & 64 & 100 & 0.01 & cosine, warmup 0.05 & \(10^{-6}\) & \(10^{-4}\) & \(x_i^{(7)}\) & \(\{0.5\}\) \\
9 & 256 & 32 & 100 & 0.01 & constant & \(10^{-6}\) & \(10^{-4}\) & \(x_i^{(9)}\) & \(\{0.25,0.5,1.0\}\) \\
\bottomrule
\end{tabular}%
}
\end{center}

For comparison, the paper-setting constant-filter baseline uses the original
paper training distribution, 10 replicas, batch size 64, 200 epochs, Adam with
learning rate 0.1, cosine scheduling with warmup fraction 0.05, and weight decay
\(10^{-6}\) for each \(N\in\{3,7,9\}\).

\section{Updated one-dimensional results}
Tables report mean \(\pm\) standard deviation over 10 independently trained
replicas.  Here \(\rNN{N}\) denotes the updated neural network,
\(\rAugConst{N}\) denotes a constant filter trained on the same augmented
sweep-selected setting as the neural network, and \(\rPaperConst{N}\) denotes the
constant filter trained with the original paper setting.  Bold entries are the
lowest mean value among the three learned-filter variants at fixed \(N\), test
case, and time.

\subsection{Gaussian}
\begin{center}
\resizebox{\textwidth}{!}{%
\begin{tabular}{c c c c c c c c c c}
\toprule
 & \multicolumn{3}{c}{\(N=3\)} & \multicolumn{3}{c}{\(N=7\)} & \multicolumn{3}{c}{\(N=9\)} \\
\cmidrule(lr){2-4} \cmidrule(lr){5-7} \cmidrule(lr){8-10}
\(t_f\) & \rNN{3} & \rAugConst{3} & \rPaperConst{3} & \rNN{7} & \rAugConst{7} & \rPaperConst{7} & \rNN{9} & \rAugConst{9} & \rPaperConst{9} \\
\midrule
0.5 & \cellbest{0.2467 \pm 0.0098} & \cell{0.9814 \pm 0.0010} & \cell{0.5772 \pm 0.0009} & \cellbest{0.1089 \pm 0.0169} & \cell{0.9618 \pm 0.0012} & \cell{0.5325 \pm 0.0012} & \cellbest{0.1108 \pm 0.0788} & \cell{0.9543 \pm 0.0042} & \cell{0.5344 \pm 0.0013} \\
1.0 & \cellbest{0.2096 \pm 0.0023} & \cell{0.9647 \pm 0.0018} & \cell{0.3636 \pm 0.0010} & \cellbest{0.1036 \pm 0.0400} & \cell{0.9237 \pm 0.0023} & \cell{0.2790 \pm 0.0013} & \cellbest{0.0363 \pm 0.0222} & \cell{0.9101 \pm 0.0081} & \cell{0.2793 \pm 0.0014} \\
\bottomrule
\end{tabular}%
}
\end{center}
\subsection{Vanishing cross section}
\begin{center}
\resizebox{\textwidth}{!}{%
\begin{tabular}{c c c c c c c c c c}
\toprule
 & \multicolumn{3}{c}{\(N=3\)} & \multicolumn{3}{c}{\(N=7\)} & \multicolumn{3}{c}{\(N=9\)} \\
\cmidrule(lr){2-4} \cmidrule(lr){5-7} \cmidrule(lr){8-10}
\(t_f\) & \rNN{3} & \rAugConst{3} & \rPaperConst{3} & \rNN{7} & \rAugConst{7} & \rPaperConst{7} & \rNN{9} & \rAugConst{9} & \rPaperConst{9} \\
\midrule
0.5 & \cellbest{0.3339 \pm 0.0091} & \cell{0.9835 \pm 0.0009} & \cell{0.6267 \pm 0.0008} & \cellbest{0.0881 \pm 0.0200} & \cell{0.9613 \pm 0.0012} & \cell{0.5252 \pm 0.0012} & \cellbest{0.1040 \pm 0.0749} & \cell{0.9552 \pm 0.0041} & \cell{0.5330 \pm 0.0014} \\
1.0 & \cellbest{0.0817 \pm 0.0048} & \cell{0.9621 \pm 0.0020} & \cell{0.2975 \pm 0.0012} & \cellbest{0.0896 \pm 0.0401} & \cell{0.9263 \pm 0.0022} & \cell{0.2841 \pm 0.0014} & \cellbest{0.0274 \pm 0.0181} & \cell{0.9106 \pm 0.0081} & \cell{0.2734 \pm 0.0015} \\
\bottomrule
\end{tabular}%
}
\end{center}
\subsection{Discontinuous cross section}
\begin{center}
\resizebox{\textwidth}{!}{%
\begin{tabular}{c c c c c c c c c c}
\toprule
 & \multicolumn{3}{c}{\(N=3\)} & \multicolumn{3}{c}{\(N=7\)} & \multicolumn{3}{c}{\(N=9\)} \\
\cmidrule(lr){2-4} \cmidrule(lr){5-7} \cmidrule(lr){8-10}
\(t_f\) & \rNN{3} & \rAugConst{3} & \rPaperConst{3} & \rNN{7} & \rAugConst{7} & \rPaperConst{7} & \rNN{9} & \rAugConst{9} & \rPaperConst{9} \\
\midrule
0.5 & \cellbest{0.3151 \pm 0.0083} & \cell{0.9829 \pm 0.0009} & \cell{0.6114 \pm 0.0008} & \cellbest{0.1050 \pm 0.0208} & \cell{0.9607 \pm 0.0012} & \cell{0.5234 \pm 0.0012} & \cellbest{0.1093 \pm 0.0977} & \cell{0.9547 \pm 0.0042} & \cell{0.5311 \pm 0.0014} \\
1.0 & \cellbest{0.1991 \pm 0.0088} & \cell{0.9623 \pm 0.0020} & \cell{0.2942 \pm 0.0012} & \cellbest{0.1479 \pm 0.0382} & \cell{0.9242 \pm 0.0023} & \cell{0.2772 \pm 0.0014} & \cellbest{0.0480 \pm 0.0227} & \cell{0.9094 \pm 0.0082} & \cell{0.2703 \pm 0.0015} \\
\bottomrule
\end{tabular}%
}
\end{center}
\subsection{Reeds}
\begin{center}
\resizebox{\textwidth}{!}{%
\begin{tabular}{c c c c c c c c c c}
\toprule
 & \multicolumn{3}{c}{\(N=3\)} & \multicolumn{3}{c}{\(N=7\)} & \multicolumn{3}{c}{\(N=9\)} \\
\cmidrule(lr){2-4} \cmidrule(lr){5-7} \cmidrule(lr){8-10}
\(t_f\) & \rNN{3} & \rAugConst{3} & \rPaperConst{3} & \rNN{7} & \rAugConst{7} & \rPaperConst{7} & \rNN{9} & \rAugConst{9} & \rPaperConst{9} \\
\midrule
5.0 & \cell{0.4493 \pm 0.0099} & \cell{0.8638 \pm 0.0066} & \cellbest{0.2613 \pm 0.0002} & \cell{0.4725 \pm 0.0813} & \cell{0.7413 \pm 0.0067} & \cellbest{0.2502 \pm 0.0002} & \cell{0.3582 \pm 0.0330} & \cell{0.7037 \pm 0.0225} & \cellbest{0.2438 \pm 0.0002} \\
10.0 & \cell{0.7526 \pm 0.0128} & \cell{0.8215 \pm 0.0077} & \cellbest{0.5550 \pm 0.0000} & \cell{0.8206 \pm 0.1134} & \cell{0.5277 \pm 0.0093} & \cellbest{0.4007 \pm 0.0006} & \cell{0.5745 \pm 0.0728} & \cell{0.4807 \pm 0.0296} & \cellbest{0.3242 \pm 0.0005} \\
\bottomrule
\end{tabular}%
}
\end{center}

\end{document}
